\documentclass[10pt,twocolumn]{article}

\usepackage[T1]{fontenc}
\usepackage{lmodern}
\usepackage[margin=0.75in]{geometry}
\usepackage{amsmath}
\usepackage{booktabs}
\usepackage{graphicx}
\usepackage[hidelinks]{hyperref}
\usepackage{xcolor}
\hypersetup{
  pdftitle={A Tracking-Free wgpu Backend: Global Pass Barriers
    Exercised in Bevy},
  pdfauthor={Dzmitry Malyshau}
}

\newcommand{\blade}{\textsc{Blade}}
\newcommand{\wgpu}{\textsc{wgpu}}

\title{A Tracking-Free wgpu Backend:\\
Global Pass Barriers Exercised in Bevy}
\author{Dzmitry Malyshau}
\date{September 2026 \\[0.4em]
\small Companion to arXiv:2607.26506. Implementation notes for JCGT;
not a rewrite of the archival study.}

\begin{document}
\maketitle

\begin{abstract}
\blade{} is an RHI that keeps Vulkan images in \texttt{GENERAL}, tracks no
per-resource state, and issues one global barrier at pass boundaries. A
cross-vendor study of that contract on matched synthetic graphs is already
archived~\cite{malyshau-sync-study}. This companion is the piece a
practitioner needs next: how to put that RHI underneath the public wgpu
API as a \emph{custom backend}, skip wgpu-core's tracker, and run a real
engine frame.

The backend translates wgpu bind groups onto \blade{} named
\texttt{ShaderData} at bind-group creation, not during
\texttt{CommandEncoder::finish} / \texttt{Queue::submit}. Bevy 0.20 then
runs unmodified except for instance construction. A headless cube+ground
cell with shadows, GPU mesh preprocessing, a depth/normal prepass, and
SSAO produces a filled 31-pass (Blade) / 46-pass (wgpu-core) frame whose
image hash agrees on a RADV Phoenix 780M. Combined record-and-submit on
that cell is \emph{not} a Blade win in the unlocked-clock pilot; the
archival claim that wgpu-core was more expensive in every synthetic cell
does not automatically transfer. Finish versus submit is backend-asymmetric:
Blade records during pass drop and submits in \texttt{Queue::submit};
wgpu-core bakes the command buffer in \texttt{finish}. Compare the sum.
MSAA resolve and swapchains remain unimplemented. The collection command
and the comparison rule are in \texttt{paper/COLLECTING.md}.
\end{abstract}

\section{Why this paper}

JCGT wants a technique one can implement, including the ugly edges
\cite{jcgt}. The archival study~\cite{malyshau-sync-study} already
answered a different question: on six GPUs, what do \emph{placement} and
\emph{pass-kind scope} of Blade's global barriers cost, relative to a
matched wgpu program, when the graph is synthetic and both sides emit the
same shaders. It did not ship a drop-in wgpu backend, did not run an
engine, and explicitly refused to attribute the host gap to tracking
alone.

This manuscript is the implementation case study. The technique is:
expose Blade as \texttt{wgpu::Instance::from\_custom}, keep the public
wgpu API, and let Bevy~\cite{bevy} drive both backends. The evaluation is
one application cell plus the same synthetic collector, on the same 780M
that the archival paper already flagged as the device where \emph{removing}
redundant barriers \emph{increased} independent-render GPU span. That is
the honest first application host: not the discrete parts that saved
$\sim$30\% on independent compute.

\section{The RHI contract, in the amount a backend needs}

Vulkan requires the application (or an RHI) to declare stage/access
dependencies and image layouts~\cite{vulkan-spec}. wgpu-core reconstructs
that state from bind groups and pass attachments~\cite{wgpu-trackers}.
NVRHI keeps per-command-list resource state even in ``manual''
mode~\cite{nvrhi-state-tracking}. Blade does none of that:

\begin{itemize}
  \item images stay in \texttt{GENERAL} for their lifetime;
  \item the encoder stores only the \emph{kinds} of passes opened since
        the last barrier (compute, render, transfer);
  \item opening a pass, or an explicit \texttt{barrier}, emits one global
        \texttt{VkMemoryBarrier} (or a pass-kind-scoped variant);
  \item the application is responsible for aliasing, exclusive
        cross-queue use, exceptional layouts, and any dependency that is
        not ``this pass kind follows that pass kind''.
\end{itemize}

The archival paper isolates two axes inside that contract
(\texttt{automatic} vs \texttt{hazard-only} placement; \texttt{Global} vs
\texttt{PassKind} scope) and compares the whole design to wgpu end to
end. Readers who need factor tables, RADV expansion of global barriers,
DCC versus FMASK under persistent \texttt{GENERAL}, or the 780M
independent-render penalty should read that study. This paper does not
repeat them.

\section{wgpu custom backend}

wgpu 30 dispatches \texttt{Instance} / \texttt{Device} / encoder traits
to wgpu-core, WebGPU, or a \emph{custom} implementation. Blade is the
custom one (\texttt{blade-wgpu}). Bevy therefore keeps a single render
path. The scientifically useful A/B is the same Bevy frame on
\texttt{--features blade} versus default wgpu-core, timed at the same
wgpu functions.

\subsection{Bind groups are not barriers}

wgpu shaders name \texttt{@group} / \texttt{@binding}. Blade shaders
export a \texttt{ShaderData} layout of named bindings; the compiler
asserts numeric bindings are absent in the native Blade path. The
backend keeps wgpu's numeric bindings, leaks stable names
\texttt{b\{binding\}} / \texttt{l\{location\}}, and maps:

\begin{itemize}
  \item \texttt{UNIFORM} $\to$ \texttt{ShaderBinding::UniformBuffer}
        (Vulkan \texttt{UNIFORM\_BUFFER}, not storage);
  \item storage buffers/textures $\to$ Blade storage;
  \item unused slots $\to$ \texttt{ShaderBinding::Unused} so wgpu's
        sparse bind groups do not shift later entries;
  \item dynamic offsets $\to$ added into \texttt{BufferPiece.offset} at
        \texttt{set\_bind\_group}.
\end{itemize}

That translation is a \texttt{profiling::scope!("blade-wgpu::create\_bind\_group")}
and a matching cost on wgpu-core's tracker side. It is \emph{not} part of
\texttt{CommandEncoder::finish} or \texttt{Queue::submit}. Host claims in
this companion always use those two entry points, via
\texttt{wgpu::util::dispatch\_stats}, so both backends are measured at the
same API. Charging bind-group construction to ``we do not track
resources'' would be a category error the archival paper already warned
against.

\subsection{Where work actually happens}

A wgpu \texttt{ComputePass} / \texttt{RenderPass} on Blade is a deferred
op list. Drop flushes into a Blade pass, which is also where the global
barrier is placed. \texttt{CommandEncoder::finish} only takes the
encoder. \texttt{Queue::submit} calls \texttt{Context::submit}, including
Blade's end-of-encoder barrier. wgpu-core does the opposite split: pass
recording fills a command buffer, \texttt{finish} bakes it,
\texttt{submit} queues it. Table~\ref{tab:split} is the comparison rule
for every host number in this paper.

\begin{table}[t]
\centering
\caption{Host split at the wgpu API. Compare the sum, not a column.}
\label{tab:split}
\small
\begin{tabular}{@{}lll@{}}
\toprule
API & Blade custom backend & wgpu-core \\
\midrule
pass Drop / end & record + global barrier & record \\
\texttt{finish} & take encoder (cheap) & bake command buffer \\
\texttt{submit} & \texttt{Context::submit} & \texttt{vkQueueSubmit} \\
bind-group create & translation, not timed here & tracker, not timed here \\
\bottomrule
\end{tabular}
\end{table}

\subsection{Limits of the scaffolding}

\texttt{blade-wgpu} is study scaffolding. It is enough to run the Bevy
cell below and the in-repo GPU tests. It is not a CTS. Deliberately
omitted, because the maximized Bevy cell does not call them, or because
they are known holes:

\begin{itemize}
  \item \texttt{create\_surface} / present --- windowed \texttt{3d\_scene}
        is not the protocol; headless offscreen is.
  \item MSAA resolve --- Blade readback of Bevy's default 4$\times$ target
        was transparent zeros. The protocol forces \texttt{Msaa::Off}.
        That is a product bug, not an environment skip.
  \item timestamp queries --- GPU elapsed is wait-to-idle
        (\texttt{Device::poll(Wait)}), recorded as
        \texttt{gpu\_timing\_method}. Do not treat it as an encoder
        timestamp.
  \item \texttt{set\_immediates}, mesh shaders, bindless, acceleration
        structures --- unimplemented; Bevy is steered off them by
        advertised features (no \texttt{INDIRECT\_FIRST\_INSTANCE} /
        \texttt{IMMEDIATES}, so mesh preprocess is
        \texttt{PreprocessingOnly} plus \texttt{NoIndirectDrawing},
        which is still compute).
\end{itemize}

Uniforms use \texttt{Queue::write\_buffer} / staging copies into
host-visible \texttt{Memory::Shared}, not push constants.

\section{Bevy cell}

Bevy constructs the wgpu instance in \texttt{initialize\_renderer}. With
\texttt{--features blade} that call is
\texttt{blade\_wgpu::BladeInstance::create\_with(\ldots).into\_wgpu()};
otherwise it is stock wgpu-core. The rest of the engine is unchanged.

The protocol example is \texttt{examples/app/sync\_bench.rs}: the same
cube and circular ground as Bevy's headless renderer, with the pass
multipliers turned \emph{on} rather than stripped for a first screenshot.

\begin{itemize}
  \item \texttt{PointLight.shadow\_maps\_enabled = true}
  \item GPU mesh preprocessing (engine default under the advertised
        limits)
  \item \texttt{DepthPrepass} + \texttt{NormalPrepass}
  \item \texttt{ScreenSpaceAmbientOcclusion} (High)
  \item \texttt{Msaa::Off}, \texttt{Tonemapping::Linear}
  \item bounded pre-roll / warmup / sample frames, CSV on stdout in the
        synthetic schema \texttt{blade-sync-bench-v1}
\end{itemize}

\texttt{paper/collect.py} launches that binary twice per repetition
(Blade feature, then wgpu-core) alongside the synthetic
\texttt{sync-bench} / \texttt{wgpu-sync-bench} matrix. The Bevy image
hash is recorded but excluded from the synthetic shader-hash agreement,
because the frame is not the matched graph. On the host below the hashes
matched anyway.

\section{Pilot on RADV Phoenix 780M}

This is machine \texttt{k6} from the archival round: AMD Radeon 780M,
RADV, Mesa 25.2.8. Clocks were \emph{not} locked
(\texttt{power\_dpm\_force\_performance\_level = auto}). The JCGT
collection used one process repetition, 8 warm-ups, 5 samples, 4
synthetic passes, $256\times256$ Bevy. That is a protocol proof and a
directional host reading, not a replacement for the ten-repetition /
1000-warmup archival matrix.

\paragraph{Correctness.}
Synthetic hashes agreed across Blade policies and wgpu-core (Khronos
synchronization validation on). The Bevy cell produced a filled
cube+ground+shadow image, 802 unique colours, hash
\texttt{fnv1a64-standard:ceae91dd19c2e122} on \emph{both} backends. A
$640\times360$ smoke (90 pre-roll frames) was the same scene, 878 unique
colours, and two consecutive Blade launches agreed. This is GPU-correct
for this scene, not a CTS.

\paragraph{Pass tally.}
\texttt{gpu\_pass\_count} is the number of wgpu
\texttt{begin\_render\_pass} + \texttt{begin\_compute\_pass} calls in a
steady sample. Blade reported 31; wgpu-core reported 46. The frontend
counter is honest on each backend; Bevy simply issues a different number
of begins once features/limits differ (timestamps, extra blits,
preprocess flavour). Do not treat $31$ versus $46$ as a barrier-count
result. Both are $\gg 1$: the cell is not a clear-only frame and not a
single blit.

\paragraph{Host cost.}
The first Bevy pilot showed Blade \texttt{submit\_ns} at $\sim$5\,ms
versus wgpu-core's $\sim$0.7\,ms. That was not \texttt{vkQueueSubmit}.
Bevy opens $\sim$28 wgpu command encoders per frame. Blade's native
loop creates one encoder and reuses it (\texttt{start} $\to$ record
$\to$ \texttt{submit}, with a wait before the next \texttt{start}).
\texttt{blade-wgpu} was instead allocating a fresh encoder --- two 1\,MiB
host-visible scratch buffers plus descriptor pools --- and
\texttt{destroy\_command\_encoder} inside \texttt{Queue::submit}.
Splitting that path: $\sim$3.5\,ms destroy, $\sim$0.66\,ms actual
queue submit (matching wgpu-core), plus $\sim$3\,ms of create that
was not even in the CSV. After pooling encoders and waiting on that
encoder's timeline before reuse, a $256\times256$ sample is:

\begin{center}
\small
\begin{tabular}{@{}lrrr@{}}
\toprule
 & record & submit & sum \\
\midrule
Blade (pooled) & $2.2$--$2.7\times 10^{4}$ & $0.57$--$0.67$\,ms & $\sim$0.6\,ms \\
wgpu-core (earlier pilot) & $0.75$\,ms & $0.67$\,ms & $1.42$\,ms \\
\bottomrule
\end{tabular}
\end{center}

The image hash is unchanged (\texttt{ceae91dd19c2e122}). Do not quote
the pre-pool 5\,ms submit as a barrier or tracking cost.

\paragraph{GPU elapsed.}
Wait-to-idle medians were $1.26$\,ms (Blade) and $1.75$\,ms (wgpu-core).
That is the same order of magnitude and is \emph{not} a timestamp-query
result. The archival 780M finding still stands as the GPU prior:
dropping fifteen automatic barriers on independent \emph{render} passes
\emph{increased} span at 32 passes, past that count's stability floor.
A Bevy frame mixes shadow, prepass, compute preprocess, SSAO, and the
main view. Nobody should expect the discrete-GPU independent-compute
saving to appear on this part until a locked-clock, timestamped
collection says so.

\section{What we will claim, and what we will not}

\textbf{Claim (implementation).}
A tracking-free RHI with global pass-boundary barriers can sit under the
wgpu API far enough for Bevy to render a multi-pass PBR frame with
shadows, a prepass, SSAO, and GPU mesh preprocessing, with an image hash
that matches wgpu-core on RADV Phoenix.

\textbf{Claim (measurement rule).}
Host A/B for that backend belongs at wgpu \texttt{finish}+\texttt{submit},
with bind-group translation reported separately, and with the finish/submit
split interpreted through Table~\ref{tab:split}.

\textbf{Claim (protocol).}
The same collector that produced the archival synthetic matrix now emits
Bevy rows in the same CSV schema. Another machine can run it; the
runbook is \texttt{paper/COLLECTING.md}.

\textbf{Not claimed.}
A GPU-time win for \texttt{GENERAL} or for dropping barriers on Bevy.
A causal ``tracking is the host cost'' decomposition. MSAA correctness.
Swapchain present. Production-quality \texttt{blade-wgpu}. Any number
from an unlocked-clock few-sample Bevy pilot as a device-span result.
The 5\,ms Blade submit from the first pilot, which was encoder
alloc/free.

\section{Limitations the next collector will hit}

Bevy on Blade still panics on several wgpu entry points. Fall back one
feature at a time (MSAA first, then contact shadows, then TAA, then
deferred lighting). Advertised limits currently omit timestamp queries
and immediates so Bevy does not call them. \texttt{max\_bind\_groups}
is 8; storage textures per stage 16, which is enough for SSAO's
plugin-load check. Cube-array textures required naga
\texttt{CUBE\_ARRAY\_TEXTURES}. Uniform uploads needed
\texttt{create\_staging\_buffer}; without it Bevy's
\texttt{ViewUniformOffset} never appeared and the 3D pass was skipped ---
a silent clear-only frame, which the protocol now rejects.

\section{Code}

\blade{} is MIT. The archival tags are \texttt{sync-study-measured-v1}
(\texttt{87ed067}) and \texttt{sync-study-v1} (\texttt{16b201e}) on
\texttt{kvark/blade}, and \texttt{blade-sync-study-v1} on
\texttt{kvark/wgpu}. This companion's sources are the \texttt{jcgt-extension} branch on
\texttt{kvark/blade} (with \texttt{blade-wgpu/}), \texttt{kvark/wgpu}, and
the \texttt{kvark/bevy} fork. Clone those three as siblings. In-repo tests
\texttt{cargo test -p blade-wgpu --test gpu} drive the shipped backend,
including a three-pass compute+prepass+color scene that asserts magenta
rather than the clear.

\bibliographystyle{plain}
\bibliography{references}

\end{document}
